\documentclass[a4paper]{article}
\usepackage[pdftex]{graphicx}
\usepackage[utf8]{inputenc}
\usepackage{enumerate}
\usepackage{siunitx}
\sisetup{locale=DE} 
\usepackage{href-ul}
\hypersetup{
	colorlinks=true,
	linkcolor=blue,
	urlcolor=blue}
\usepackage{icomma}
\usepackage{amssymb}
\usepackage{geometry}
\geometry{a4paper, top=15mm, left=15mm, right=15mm, bottom=15mm,
	headsep=10mm, footskip=12mm}

\begin{document}
{\bf Wärmelehre}
\begin{enumerate}[1.]
{\bf \item ,,Unter Wärme versteht man die Energie, die infolge von Temperaturunterschieden zwischen Körpern übertragen wird"} --- Erkläre kurz die drei Möglichkeiten der Energieübertragung!


%Aufgabe 2
\begin{minipage}{0.5\textwidth}
{\bf \item Im Diagramm rechts ist ein qualitativer Tempe\-raturverlauf eines Vorgangs dargestellt,} bei dem einem Stoff kontinuierlich Wärme zugeführt wurde.
\begin{enumerate}[a)]
\item Um welchen Stoff könnte es sich handeln? 
\item Erläutere die Begriffe Schmelzen, Erstarren, Sieden, Kondensieren, Sublimieren und Resublimieren.
\item Erläutere was in den Kurvenpunkten A bis E geschieht 
\end{enumerate}
\end{minipage}
\hfill
\begin{minipage}{0.45\textwidth}
\includegraphics[width=5 cm]{wasser105.pdf}
\end{minipage}

{\bf \item  160 g Eis von $\SI{0}{\celsius}$ stehen in einem Gefäß auf einer Heizplatte.} Das Eis soll vollständig in Dampf von $\SI{100}{\celsius}$ umgewandelt werden. [$r_{H_2O} = \SI{2257}{\kilo\joule\over \kilogram};\quad s_{H_2O}=\SI{334}{\kilo\joule\over \kilogram};\quad c_{H_2O}=\SI{4,18}{\kilo\joule\over \kilogram\kelvin}$]
\begin{enumerate}[a)]
\item Welche Wärmemenge muss dem Eis zugeführt werden?
\item Wie lange muss die Heizplatte die entsprechende Wärme abgeben, wenn sie eine  Heizleistung von $P=\SI{2800}{\watt}$ hat?
\end{enumerate}


\begin{minipage}{0.5\textwidth}
{\bf \item Nebenstehender Schaltplan zeigt einen Stromkreis mit einer Spannungsquelle, einem Verbraucher und einem Bimetallstreifen.}
\begin{enumerate} 
\item Erläutere die Funktionsweise eines Bimetallstreifens.
\item Benenne diese Schaltung und gib eine Anwendung dieser an.  
\end{enumerate} 
\end{minipage}
\hfill
\begin{minipage}{0.45\textwidth}
\includegraphics[width=6 cm]{bimetall105.pdf}
\end{minipage}



{\bf \item Sauerstoff, der mit einem Anteil von ungefähr 21 \% in der Luft vorhanden ist, verflüssigt sich bei etwa $ T=\SI{90}{\kelvin}$ über dem absoluten Temperaturnullpunkt.}\\
Um wie viel Kelvin muss man den Sauerstoff der Luft bei einer Zimmertemperatur von
 $\SI{19}{\celsius}$  abkühlen, damit er flüssig wird ? 


{\bf \item $m_1=\SI{200}{\gram}$ Wasserdampf der Temperatur $\vartheta_{d} =\SI{100}{\celsius}$ werden in $m_2=\SI{1800}{\gram}$ Wasser der Temperatur $\vartheta_{w} = \SI{20}{\celsius} $ so eingeleitet, dass der Wasserdampf vollständig kondensiert.} Die Temperatur der Flüssigkeit ist dann 
$\vartheta_{m} = \SI{82}{\celsius} $. Der Energieaustausch mit der Umgebung und dem Kalorimetergefäß soll hier vernachlässigt werden.
\begin{enumerate}[a)]
\item Welche Wärmeenergie nimmt das Wasser der Masse $m_2$ auf?
\item Berechne die spezifische Kondensationsenergie des Wasserdampfs!
\end{enumerate}

{\bf \item $m_w=3,2~kg$ Wasser der Temperatur $\vartheta_w = \SI{30}{\celsius}$ und $\SI{1,5}{\kilogram}$ Eis werden in einem Kalorimeter gemischt. Nachdem das Eis zu 70\% geschmolzen ist, ändert sich nichts mehr.} Welche Temperatur hatte das Eis am Anfang?
[Spezifische Wärmekapazität von Eis  $c_{eis}= \SI{2,1}{\kilo\joule \over \kilogram\kelvin}$]


%Aufgabe 1
{\bf \item Stellt man einen langen Metalllöffel in eine heiße Tasse Kaffee, so wird auch der aus dem Kaffee herausragende Teil des Löffels heiß.}
\begin{enumerate}[a)]
\item  Erkläre mit dem Teilchenmodell, wie diese Wärmeleitung funktioniert.
\item Erläutere, wie die Übertragung innerer Energie durch Konvektion vor sich geht. Erläutere dies dann auch an einem selbst gewählten Beispiel.
\end{enumerate}

%Aufgabe 2
{\bf \item Es soll berechnet werden, wie sehr sich die Trommelbremse eines Fahrrades aus $\SI{1,0}{\kilogram}$  Stahl erwärmt,} wenn während einer Bergabfahrt mit einem Höhenunterschied von 100~m ständig gebremst wird. Die spezifische Wärmekapazität von Stahl ist $c=\SI{0,48}{\kilo\joule \over \kilogram\kelvin} $, die Masse des Fahrrades mit Fahrer $\SI{85}{\kilogram}$.

{\bf \item Wie lautet der erste Hauptsatz der Wärmelehre?}

\fbox{
	\begin{minipage}{0.5\textwidth}
		Zur Lösung bitte \href{https://www.okuyakl.de/phys/p9thaL402/ll402.pdf}{hier klicken} oder den QR-Code scannen.\\
	Weitere Arbeitsblätter gibt es unter 
	
	\href{https://www.okuyakl.de}{www.okuyakl.de}
	\end{minipage}
	\hfill
	\begin{minipage}{0.4\textwidth}
		\includegraphics[width=1.5 cm]{../../viecher/zwe03}
		\includegraphics[width=3 cm]{qr402}
		\includegraphics[width=2 cm]{../../viecher/afanticon1}
		
	\end{minipage}}

\end{enumerate} 
\end{document}%Lösung-------------------------------------------
